\subsection{Fibrewise smoothness can improve at specialization}

\paragraph{Theorem E4.S.}
There is a projective flat family of geometrically integral reduced surfaces
over $R=\overline{\mathbf F}_2[[t]]$ whose geometric generic Picard
tau-part is nonsmooth, whereas its special Picard tau-part is smooth.
Thus the locus of fibrewise smoothness is not open. This is a precise
negative answer to the smoothness interpretation of 57:309(b$'$);
it does not identify fibrewise smoothness with relative smoothness.

Let $C\subset\mathbb P^2_R$ be the cubic
\[
 y^2z+txyz+tyz^2=x^3,\qquad o=[0:1:0].                 \tag{E4.S.1}
\]
Its discriminant and $c_4$ are
$\Delta=t^4(1+t^2)$ and $c_4=t^4$. Hence its generic fibre is a
smooth ordinary elliptic curve, with $j=t^8/(1+t^2)\ne0$.
The special fibre is the integral cuspidal cubic $y^2z=x^3$.
The coefficient of $y^2z$ makes $o$ smooth over $R$.
The involution
\[
 \sigma[x:y:z]=[x:y+tx+tz:z]                           \tag{E4.S.2}
\]
is inversion generically and fixes $o$.
Let $B/R$ be the constant ordinary elliptic curve
$y^2z+xyz=x^3+z^3$, let $P=(0,1)$ on its affine chart,
and let $\Gamma=(\mathbb Z/2\mathbb Z)_R$ act on $C\times_R B$
through $(\sigma,t_P)$. The nonzero point $P$ has order two, so this
action is free. Set
\[
 X=(C\times_R B)/\Gamma,\qquad Y=B/\langle P\rangle.
\]
The finite-group quotient is projective: an orbit tensor product of
an ample line bundle supplies an ample linearized line bundle.
Its quotient map is finite \'etale, so flatness descends, and all
geometric fibres of $X/R$ are integral and reduced. The invariant
section $o\times B$ gives a section $i:Y\to X$.

Write $G=\operatorname{Pic}^0_{C/R}$. It is a smooth separated group
scheme of finite presentation: the projective integral curve family
has its Picard scheme, and $H^2(C_s,\mathcal O_{C_s})=0$ removes
all infinitesimal obstructions. The degree-zero locus is its neutral
part and also its tau-part. The equality $\sigma^*=[-1]$ on $G$
holds generically; it holds everywhere because $G$ is flat over
the trait and the target is separated.
There is a base-change compatible isomorphism
\[
 \operatorname{Pic}^\tau_{X/R}
       \simeq Y^\vee\times_R G[2].                  \tag{E4.S.3}
\]
Here is the descent argument, including the singular special fibre.
Rigidification along $o$ makes a representing line bundle on $C$
unique up to a unique rigidified isomorphism. For a point of $G[2]$
the unique isomorphism with its $\sigma$-pullback respects this
rigidification, and its square is the identity. It therefore supplies
a canonical $\Gamma$-linearization. Pulling to $C\times B$ and
descending gives a line bundle on $X$ trivial on $i(Y)$.
Together with pullback from $Y$, this defines (E4.S.3).
Conversely, subtract the restriction to $i(Y)$ from a tau-trivial
line bundle on $X$. Its pullback to $C\times B$ is a bundle from
$C$ alone, by the product formula for tau and its trivialization on
$o\times B$. Descent makes its class $\sigma$-invariant. Its
linearization agrees with the canonical one: their ratio is a unit
on the proper connected product and is one along $o\times B$.
These constructions are inverse and commute with arbitrary base
change. Numerical triviality is detected by the finite surjective
cover, by the degree formula on curves. This proves (E4.S.3) as a
group-scheme assertion.

Over a geometric generic point,
$G[2]\simeq\mu_2\times(\mathbb Z/2\mathbb Z)$, so (E4.S.3)
is nonsmooth and has dimension one. On the special fibre,
$G_s\simeq\mathbb G_a$. Indeed, the cusp has normalization
$x=u^2,y=u^3$, conductor $(u^2)$, and the units quotient
\[
 (k[u]/(u^2))^\times/k^\times\simeq\mathbb G_a.
\]
The normalization has no degree-zero Picard group, giving the
assertion. Since $[2]$ is zero on $\mathbb G_a$,
\[
 (\operatorname{Pic}^\tau_{X/R})_s
       \simeq Y_s^\vee\times\mathbb G_a
\]
is smooth of dimension two. Its tangent dimension is two, as is
the generic tangent dimension, whereas the defect
$\dim\operatorname{Lie}-\dim$ drops from one to zero.
Consequently that defect is not upper semicontinuous either.
Relative smoothness remains an open condition; in this example it
fails at every special point because the relative dimension jumps.
The separate dimension function is upper semicontinuous, in agreement
with FGA VI Corollary~2.7, p.234.

The finite polynomial identities, Jacobian checks, discriminant,
and cusp parameterization are recorded in
\texttt{code/check\_cuspidal\_igusa.output.json}.
The singular-cubic extension of the rigidified descent proof received
a separate targeted AI audit from /root/dgj, which found no gap.
